% ══════════════════════════════════════════════════════════════════════════════
\chapter{Results and Validation}
\label{ch:results}
% ══════════════════════════════════════════════════════════════════════════════

\section{Slew Speed}
\label{sec:slew-speed}

Slew speed was measured on a single axis, unloaded, at 24~V supply, by commanding the rotator to traverse 360\degree\ at a fixed PWM duty and timing the traverse with the onboard encoder tick counter. Four duty levels were tested. Raw data and per-run encoder-derived rates are listed in Appendix~\ref{app:pid}.

\begin{table}[H]
    \centering
    \begin{tabular}{lll}
        \toprule
        Duty (\%) & Speed (\degree/s) & 360\degree\ time (s) \\
        \midrule
        35 & 11.2 & 32.1 \\
        55 & 18.9 & 19.0 \\
        75 & 26.9 & 13.4 \\
        95 & 35.0 & 10.3 \\
        \bottomrule
    \end{tabular}
    \caption{Slew-speed measurements, unloaded, 24~V supply.}
\end{table}

Linear regression over the four measured points gives slope 0.40\degree/s per \% duty, $R^2 > 0.99$. 100\,\% duty was tested separately and trips the TB6642FG overcurrent protection on motor stall, so it is not included in the fit. The 35\degree/s ceiling at 95\,\% duty exceeds the total angular rate $\dot{\theta}_{\max} \approx 1.1$\degree/s derived in \S\ref{sec:orbital-mechanics} by a factor of approximately 32. It does not cover the instantaneous \textit{azimuth} rate during near-zenith passes (which can exceed 10\degree/s for a few seconds); \S\ref{sec:tracking-accuracy} reports the resulting AZ tracking error during such passes.

\section{PID Step Response}
\label{sec:pid-step}

With the deployed gains ($K_p = 0.15$, $K_i = 0.03$, $K_d = 0.02$, duty-filter $\alpha = 0.15$, max duty 0.95; \S\ref{sec:rotator-pid}) a 90\degree\ commanded step settles to within the 0.05\degree\ deadband in approximately 4~s, with overshoot below 0.5\degree. The settling time decomposes as: slew-limited traverse of 90\degree\ at 35\degree/s peak rate = $90/35 = 2.6$~s, plus roughly 1.4~s of final convergence through the duty output filter and deadband. The closed-loop bandwidth is bounded above by the duty-output filter time constant:
\begin{equation}
    \tau_\text{duty} = \frac{\Delta t}{\alpha} = \frac{0.01}{0.15} \approx 0.067~\text{s}, \qquad f_{-3\,\text{dB}} \approx \frac{1}{2\pi \tau} \approx 2.4~\text{Hz}.
\end{equation}

In practice the slew-rate limit sets the large-signal response time; the 2.4~Hz figure applies to small-signal tracking above the deadband. Either is well above the $\leq$1.1\degree/s total angular rate of a 400~km LEO pass (\S\ref{sec:orbital-mechanics}).

\begin{figure}[H]
    \centering
    \includegraphics[width=0.8\textwidth]{figures/pid_step_overlay.png}
    \caption{PID step response for a 90\degree\ commanded step.}
\end{figure}

\section{Satellite Tracking Accuracy}
\label{sec:tracking-accuracy}

Tracking error was computed per pass from the CSV log (commanded AZ/EL from \code{station.py} pass predictor vs. actual encoder readback from the MotorPCB, sampled at 5~Hz; logs in \file{captures/20260409/pass\_logs/} and \file{captures/20260412/}). Absolute-value errors are binned by the satellite's peak-elevation category:

\begin{table}[H]
    \centering
    \begin{tabularx}{\textwidth}{lX}
        \toprule
        Scenario & Error ($|$commanded $-$ actual$|$) \\
        \midrule
        EL tracking, all passes & $<$0.5\degree\ throughout the pass \\
        AZ tracking, 20\degree--70\degree\ peak elevation & 0\degree--3\degree\ throughout the pass \\
        AZ tracking, peak elevation $>$70\degree\ (near-zenith) & 11\degree--166\degree\ during the $\sim$2--4~s zenith-crossing window; recovers to 0\degree--3\degree\ below 70\degree \\
        \bottomrule
    \end{tabularx}
    \caption{Tracking error binned by satellite peak elevation.}
\end{table}

Across passes with peak elevation in the 20\degree--70\degree\ band (the range where the Chapter~\ref{ch:background} link budget closes), both axes remain inside the Siretta Oscar 44 3-dB beamwidths (54\degree\ H-plane, 48\degree\ E-plane; \S\ref{sec:problem-statement}). The near-zenith AZ excursions exceed the 35\degree/s motor ceiling from \S\ref{sec:slew-speed} for a few seconds per pass; this is a speed-limit event, not an accuracy-limit one, and the error decreases monotonically as the satellite descends below 70\degree\ elevation.

\section{RF Devboard Validation}
\label{sec:rf-measurements}

Two devboards were connected to a host PC and exercised with \file{cc1200\_rf\_test.py}. The test script loads a bench validation SmartRF profile (2.4\,kbps 2-GFSK, 5.2\,kHz deviation, 10\,kHz channel filter; modulation index $h = 2\Delta f/f_b \approx 4.3$, chosen for robustness at short range rather than spectral efficiency, this is not the AetherSpace flight configuration, which remains to be defined once the satellite's modulation parameters are finalised), tunes both boards to a specified frequency, and runs a bidirectional packet exchange: Board\,A transmits a sequence of fixed-size payloads while Board\,B receives, then the roles are reversed. Payload matching is verified byte-for-byte.

The frequency sweep covered 430--439\,MHz in 1\,MHz steps with 50 packets per direction per frequency (1000 total). Test conditions: 1\,m board separation, dipole antennas, indoor lab, CRC-16 enabled. Table~\ref{tab:devboard-bench} shows the per-frequency results averaged over both directions.

\begin{table}[H]
    \centering
    \caption{Devboard bench test results per frequency, averaged over both directions (50 packets each). LQI lower~=~better.}
    \small
    \begin{tabular}{rrrrr}
        \toprule
        Freq (MHz) & RSSI avg (dBm) & LQI avg & Latency avg (ms) & Loss (\%) \\
        \midrule
        430 & $-$83.6$^\dagger$  & 0.4 & 103.5 & 0.0 \\
        431 & $-$103.7           & 0.3 & 103.3 & 0.0 \\
        432 & $-$102.4           & 0.5 & 103.3 & 0.0 \\
        433 & $-$103.6           & 0.4 & 104.1 & 0.0 \\
        434 & $-$103.7           & 0.4 & 104.0 & 0.0 \\
        435 & $-$88.8$^\dagger$  & 0.5 & 106.5 & 0.0 \\
        436 & $-$103.2           & 0.4 & 104.8 & 0.0 \\
        437 & $-$103.9           & 0.5 & 103.5 & 0.0 \\
        438 & $-$103.9           & 0.4 & 104.3 & 0.0 \\
        439 & $-$103.7           & 0.4 & 103.6 & 0.0 \\
        \bottomrule
        \multicolumn{5}{l}{\footnotesize $^\dagger$ Fractional-N PLL low-dither artefact; see text.}
    \end{tabular}
    \label{tab:devboard-bench}
\end{table}

Zero packet loss and 100\,\% payload match across all 10 frequencies confirm that the CC1200 SPI driver, SmartRF profile loading, COBS framing, TX and RX FIFO paths, sync-word detection, CRC validation, and bidirectional TRX switching all function correctly on the devboard hardware. LQI is a CC1200 demodulator quality metric that accumulates constellation error over 64 symbols after sync detection; lower values indicate smaller error. The uniform averages of 0.3--0.5 across all frequencies reflect near-perfect symbol recovery at short range, and confirm that neither the RSSI anomalies at 430 and 435\,MHz nor the noise-floor readings elsewhere affect demodulator performance. Latency averages of 103--107\,ms are consistent across the sweep, reflecting the end-to-end pipeline from host TX command through USB CDC, CC1200 transmission, and USB CDC return; the flat profile confirms no frequency-dependent timing effects.

The elevated RSSI at 430 and 435\,MHz is a measurement artefact rooted in the CC1200's fractional-N PLL. In a fractional-N synthesiser the divider ratio is not an exact integer, so a sigma-delta modulator continuously switches the divider between neighbouring integer values, a process called dithering, which spreads quantisation noise across a broad frequency range and produces typical phase-noise behaviour. At 430\,MHz the FREQ register evaluates to an exact integer multiple (\textsc{freq2}\,=\,\texttt{0x2B}, \textsc{freq1}\,=\,\textsc{freq0}\,=\,\texttt{0x00}), placing the synthesiser at a true integer-N boundary where the sigma-delta is idle and dithering does not occur. At 435\,MHz the fractional part is exactly one-half (\textsc{freq2}\,=\,\texttt{0x2B}, \textsc{freq1}\,=\,\texttt{0x80}, \textsc{freq0}\,=\,\texttt{0x00}), so the sigma-delta alternates between only two adjacent divider states rather than spanning a broader range; dithering is minimal rather than absent. Both conditions produce lower-than-typical local-oscillator phase noise, which, combined with the RSSI sample being taken immediately after packet reception before the running average has decayed to the ambient floor, produces the apparent peaks.

What the bench test does not establish is satellite reception, the bench test scopes the transceiver stack in isolation, not the system sensitivity in a satellite geometry.

\section{Field Deployment and Live-Pass Observations}
\label{sec:field-tests}

Three field sessions were conducted with the complete station on a camera tripod in an open field, powered from a bench supply via extension cord and using an iPhone hotspot for internet. Total setup time from unpacking to first tracked pass was approximately 4 minutes, with no SSH session or laptop-side configuration required beyond opening the dashboard in a browser.

\begin{figure}[H]
    \centering
    \includegraphics[width=0.75\textwidth]{photos/station_wide.jpg}
    \caption{Wide shot of the deployed station in the field.}
\end{figure}

CSV pass logs (timestamp, commanded AZ/EL, encoder AZ/EL, RSSI, frequency, Doppler offset, slant range, packet count, 5~Hz cadence) were written to \file{captures/20260409/pass\_logs/} and \file{captures/20260412/}. Across all three sessions 210 passes were logged. The Doppler correction loop ran throughout: the observed $\pm$9~kHz sweep at 432~MHz is consistent with the $\pm$11~kHz maximum predicted in \S\ref{sec:doppler}, with the commanded offset clipped at the peak radial velocity the satellite actually reaches. The noise floor and reception outcome derived from these logs are reported in \S\ref{sec:packet-reception} and are not repeated here.

\subsection{Firmware bugs surfaced by live passes}

Four firmware/software defects were diagnosed during field sessions and fixed in the same session:

\begin{itemize}
    \item AZ integral windup after zenith crossing. When the AZ axis saturated during a high-rate zenith crossing, the PID integral accumulated a large error; after the satellite descended and the rate dropped, the integral fought the proportional term and oscillated. Fixed with the saturation-based anti-windup described in \S\ref{sec:rotator-pid}: the integral stops accumulating when $|u| \geq 0.95$ and the error sign would grow it further.
    \item Duty output filter too slow. The PWM output LPF at $\alpha = 0.05$ (time constant 0.2~s) made the AZ axis unresponsive at tracking rates below $\sim$1\degree/s. $\alpha$ was raised to 0.15 (0.067~s) with no observed increase in motor-current transients; this is the deployed setting used throughout \S\ref{sec:pid-step}.
    \item Doppler base-frequency compounding. A defect in \code{set\_frequency()} overwrote the per-satellite base frequency with the Doppler-corrected value, so each successive 2~Hz update accumulated on top of the last correction. Fixed by introducing a separate \code{uhf\_base\_freq\_hz} variable that holds the satellite's nominal frequency constant across the pass, with the Doppler offset applied only when writing the CC1200 FREQ/FREQOFF registers.
    \item Rotator socket timeout. The TCP \code{set\_position()} call to \code{rotctld} blocked up to 5~s on the receive side, limiting effective tracking rate. Reduced to 0.5~s with fire-and-forget sends, plus a 15~s keepalive ping during AOS wait to keep the socket responsive. Implementation in \code{station.py}.
\end{itemize}

\section{Packet Reception: Result and Interpretation}
\label{sec:packet-reception}

As of April 2026, no protocol-valid frames have been decoded from orbit. Across the three field sessions the station tracked approximately 210 satellite passes (16 on April 1 baseline with uncalibrated RSSI readout, 184 on April 9, and 10 on April 12) and captured tens of megabytes of raw binary data, but:

No CRC-16/CCITT-valid HDLC frames (with or without G3RUH descrambling), no AX.25 frames with valid FCS, no AX.100 Mode~5 frames passing Reed--Solomon RS(255,223) correction, and consequently no decoded packets submitted to the SatNOGS DB from real satellites.

The receive chain has been validated for local test frames in a short-range ground-to-ground configuration between two devboards (\S\ref{sec:rf-measurements}): transmitted test frames decode at the receiver with valid sync word and zero payload errors across all tested frequencies, exercising the RF path, SmartRF register set, FIFO drain, and host-side decoders under those conditions. This ground validation confirms the digital and RF path is internally consistent; it does not establish that the system is receiving orbital signals.

The April 1 session predates the RSSI-readout fix (AGC\_GAIN\_ADJUST calibration) and is used as the negative control in the analysis below; the April 9 and April 12 sessions use the post-fix calibration.

\subsection{Reception Modes Deployed}

Two RX modes were used, selected per satellite by \code{station.py} (\S\ref{sec:rf-hat-firmware}):

\begin{enumerate}
    \item FIFO mode (default): CC1200 matches a sync word (e.g. AX.100 ASM \code{0x930B51DE}), then delivers fixed-length 255-byte chunks. If sync or modulation parameters don't match, the FIFO drains receiver noise after false sync triggers.
    \item Serial mode with AX.25 G3RUH decode: for $\geq$4800~bps FSK, the firmware switches to PIO serial capture and the G3RUH / NRZ-I / HDLC pipeline of \S\ref{sec:serial-mode}. No CRC-valid AX.25 frames were recovered from orbital data.
\end{enumerate}

\subsection{Why no decoded frames}

The \S\ref{sec:link-budget} link budget shows the AetherSpace downlink is already marginal at the datasheet CC1200 sensitivity ($-$5.7~dB theoretical margin at 30\degree\ elevation). In the field, board-level interference raises the effective noise floor to $-$96~dBm, burying the received signal of $-$125.7~dBm approximately 30~dB below the detection threshold. CRC-16 is error detection, not correction; RS(255,223) cannot compensate for a deficit of that size. The zero-frame outcome is the predicted outcome.

Three factors contribute:

\begin{enumerate}
    \item \textbf{Front-end noise figure and board interference.} No LNA, so the CC1200 AGC operates on antenna noise plus on-board interference (LMR51420 buck harmonics, RP2040 digital, Pi GPIO/HDMI); measured noise floor sits approximately 30~dB above the expected kTB+NF limit at the CC1200 input.
    \item \textbf{Modulation parameter uncertainty.} The SatNOGS DB does not always publish exact symbol rate / deviation / sync word for non-AetherSpace satellites; any mismatch in FIFO mode produces zero output.
    \item \textbf{Link geometry.} A +14~dBm / +2~dBi CubeSat leaves negative theoretical margin at mid-elevation even with a perfect receiver.
\end{enumerate}

\subsection{What was observed: suggestive but not diagnostic}

The captures contain statistical patterns that are individually compatible with weak sub-threshold satellite reception, but each is also compatible with receiver noise or demodulator artefacts triggered by RF activity at the sync-word matcher. Both hypotheses are named explicitly below.

\paragraph{A. Elevation-binned mean RSSI (suggestive).} Grouping 62 tracked passes by peak elevation, mean peak RSSI increases monotonically by +4.2~dB from the 0--20\degree\ bin to the 60--90\degree\ bin. Weak satellite RF entering the front end at higher elevations is one possible explanation. However, the dataset is also confounded by pointing-dependent antenna temperature, local 433~MHz interference, AGC behaviour, and pass-selection effects; a passive antenna pointed at cold sky would normally lower the noise floor relative to the horizon, so this trend cannot be treated as proof of satellite reception without calibrated controls.

\paragraph{B. Cross-capture byte matching on SNUGLITE-II (suggestive, inconclusive).} Four independent SNUGLITE-II captures contain shared subsequences up to 81 bytes long. Two hypotheses are compatible: (i) repeated beacon frames captured at different orbital points, and (ii) a CC1200 demodulator/idle pattern emitted repeatedly at false sync triggers. The April 9 log explicitly identifies a 99-byte CC1200 idle pattern appearing identically across all satellites at all frequencies; the SNUGLITE-II matches have not yet been shown to differ structurally from this idle pattern, so (ii) is not excluded.

\paragraph{C. Pass-profile RSSI range (weak).} A small subset of passes show an AOS-to-peak-to-LOS RSSI excursion of several dB, loosely consistent with $1/r^2$ path-loss modulation of a real signal but equally consistent with antenna noise-temperature variation as the beam sweeps across ground and sky. Not diagnostic.

\begin{figure}[H]
    \centering
    \includegraphics[width=0.8\textwidth]{figures/evidence/10_sky_coverage.png}
    \caption{Sky coverage heatmap: RSSI over AZ/EL across all tracked passes.}
\end{figure}

\paragraph{D. Doppler-sweep profile (confirms orbital geometry, not reception).} The expected $\pm$9~kHz sweep at 432~MHz with zero-crossing at TCA is observed on every pass. Because this sweep is \textit{applied} by station software to the FREQOFF register regardless of whether any signal is present, it verifies correct predict/correct software behaviour only.

\begin{figure}[H]
    \centering
    \includegraphics[width=0.8\textwidth]{figures/evidence/06_doppler_scurve.png}
    \caption{Doppler-correction S-curve applied during one pass.}
\end{figure}

\paragraph{E. Capture-byte entropy (weak).} Captures show byte entropy of 6.95--7.91 bits/byte, below the 8.0 bits/byte expected of uniform random noise. Filtered or correlated noise from a demodulator can produce similar signatures without encoding any information. The entropy is consistent with ``not white noise'' but not with ``structured satellite telemetry'' specifically.

\paragraph{F. MIMAN AX.100 Mode~5 analysis (inconclusive).} MIMAN captures processed through AX.100 Mode~5 achieved Golay(24,12) decode rates of 73\,\% (April 9) and 66\,\% (April 12). The random false-positive baseline for Golay(24,12) with 3-error correction is approximately 57\,\%, so the observed rates are only marginally above chance. Zero frames passed the subsequent RS(255,223) correction. Golay matches without downstream RS-decoded content are not a definitive indicator. See \file{captures/evidence/miman\_decode\_report.json}.

\subsection{Academic Conclusion}

The primary finding is established independently of observations A--F: the theoretical margin at 30\degree\ elevation is already $-$5.7~dB at datasheet CC1200 sensitivity, and the field-measured noise floor of $-$96~dBm buries the $-$125.7~dBm received signal by approximately 30~dB, making protocol-valid decode physically impossible regardless of which secondary hypothesis explains the statistical patterns. The secondary finding is that A--F are individually consistent with both weak sub-threshold reception and receiver-internal noise artefacts, and cannot be used to distinguish between the two hypotheses. This ambiguity is expected and does not alter the root-cause diagnosis: the board-noise floor sits approximately 30~dB above the thermal\,+\,NF baseline. A mast-mounted front-end LNA is the primary intervention: its gain lifts the received signal above the fixed board-interference floor before it reaches the CC1200. However, with the field noise floor at $-$96~dBm, a 20~dB LNA still leaves the signal approximately 10~dB below the detection threshold, so the LNA must be combined with board-level EMI mitigation: improved buck-converter output filtering, a post-buck LDO on the CC1200 supply rails to suppress conducted switching noise, and a SAW bandpass filter on the antenna port to reject out-of-band interference before it reaches the LNA (\S\ref{sec:future-work}).

\subsection{AetherSpace-Specific Outlook}

For AetherSpace the modulation-match risk is lower than for arbitrary satellites because the CubeSat and the ground station use the same CC1200 transceiver, provided both sides load the same finalised SmartRF register set. The remaining risk is sensitivity, addressed in \S\ref{sec:future-work}. The software architecture already exposes an \code{RfBackend} abstraction with both CC1200 and RTL-SDR implementations, allowing an SDR-based wideband IQ capture to serve as an independent verification path once an LNA is installed.

\section{Cost Comparison}

The v2 MOTOR\_RF\_HAT BOM (from \S\ref{sec:cost}) compared against commercial and existing open-source alternatives. The ``rotator'' column counts the controller board + motors + mechanics only, excluding the station computer, antenna, coax, power supply, and tripod; ``total'' adds a Pi 3A+.

\begin{table}[H]
    \centering
    \small
    \begin{tabular}{@{}llll@{}}
        \toprule
        System & Type & Rotator (\euro{}) & +Pi (\euro{}) \\
        \midrule
        Yaesu G-5500   & Commercial    & 760--900   & 1200+ \\
        SPID RAS       & Commercial    & 1200--1300 & 1800+ \\
        SatNOGS v3     & Stepper, belt & 300--500   & 600+ \\
        \textbf{This work (v2)} & \textbf{DC + gears} & \textbf{$\sim$75} & \textbf{$\sim$100} \\
        \bottomrule
    \end{tabular}
    \caption{Cost comparison. v2 is design-complete but not yet fabricated; $\sim$\euro{}75/\euro{}100 are BOM projections (LCSC + JLCPCB, April 2026).}
\end{table}
